% =============================================================================
%  Thesis Pipeline — Computational Summary
%  Auto-generated from src/ code analysis
% =============================================================================
\documentclass[11pt,a4paper]{article}

% ── packages ─────────────────────────────────────────────────────────────────
\usepackage[utf8]{inputenc}
\usepackage[english]{babel}
\usepackage[margin=1in]{geometry}
\usepackage{amsmath,amssymb,amsfonts}
\usepackage{booktabs}
\usepackage{longtable}
\usepackage{array}
\usepackage{multirow}
\usepackage{xcolor}
\usepackage{colortbl}
\usepackage{enumitem}
\usepackage{hyperref}
\usepackage{cleveref}
\usepackage{caption}
\usepackage{graphicx}
\usepackage{tabularx}
\usepackage{pdflscape}

\hypersetup{
    colorlinks=true,
    linkcolor=blue!60!black,
    citecolor=green!50!black,
    urlcolor=blue!70!black,
}

\newcommand{\code}[1]{\texttt{#1}}
\newcommand{\var}[1]{\textit{#1}}
\newcommand{\mat}[1]{\mathbf{#1}}
\newcommand{\R}{\mathbb{R}}
\newcommand{\PS}{\mathrm{PS}}
\newcommand{\ZCI}{\mathrm{ZCI}}
\newcommand{\BC}{\mathrm{BC}}

\title{Computational Pipeline Summary:\\
       Zoobenthic Community Index Construction\\[6pt]
       \large Methods, Data Flow, and Quality Assessment}
\author{}
\date{\today}

\begin{document}
\maketitle
\tableofcontents
\clearpage

% ═════════════════════════════════════════════════════════════════════════════
\section{Overview}
\label{sec:overview}
% ═════════════════════════════════════════════════════════════════════════════

The pipeline constructs a \emph{Zoobenthic Community Index} (ZCI) that
quantifies benthic community condition along a sediment-contamination
gradient in the St.~Clair--Detroit River System (SCDRS) corridor.
The corridor-wide workflow is orchestrated by
\code{run\_SCDRS\_full\_pipeline.py} and proceeds through the following stages:

\begin{enumerate}
    \item \textbf{Stage~1} — Pollution PCA: derive a composite \emph{Pollution Score}
          from 16 sediment-chemistry variables.
    \item \textbf{Stage~RDA} — Redundancy Analysis (complementary validation):
          quantify how much of the community variation is explained by
          environmental predictors.
    \item \textbf{Stage~2} — Taxa assemblage clustering: partition
          least-polluted (reference) sites into habitat groups using
          Ward's hierarchical clustering on taxa data.
    \item \textbf{Hindsight relabel + ANOVA}: merge small clusters and
          test environmental/taxa differences across habitat groups.
    \item \textbf{Stage~3} — LDA classification: train a Linear
          Discriminant Analysis on reference sites to classify all
          remaining (non-reference) sites into habitat groups.
    \item \textbf{Stage~4} — Bray--Curtis NMDS + ZCI: ordinate taxa
          communities by non-metric multidimensional scaling (NMDS) on
          Bray--Curtis dissimilarities, then compute a ZCI for every site.
    \item \textbf{Stage~5} — Piecewise quantile regression: model the
          ZCI--Pollution relationship with piecewise linear quantile
          regression and wild-bootstrap confidence intervals.
\end{enumerate}

Stages~1--4 and the hindsight step form the \textbf{core analytical
stream}; Stage~RDA is \textbf{complementary/validation}; Stage~5 is a
\textbf{post-hoc inference} layer.  A prospective \textbf{TITAN}
(Threshold Indicator Taxa ANalysis) module is discussed in
\cref{sec:titan}.

Each stage reads one or more \emph{artifact} Excel workbooks produced by
its predecessors, adds new columns, and writes an updated artifact.
All workbooks use a three-level \code{MultiIndex} column header
\code{(block, level, variable)}.


% ═════════════════════════════════════════════════════════════════════════════
\section{Data Dictionary and Notation}
\label{sec:data}
% ═════════════════════════════════════════════════════════════════════════════

\subsection{Workbook Structure}
\label{sec:data:workbook}

The master workbook
(\code{complete\_env\_taxa\_chemical\_Feb\_3.xlsx}) is read by
\code{read\_study\_data()} (module \code{zci.io.readers}) into a
\code{pandas.DataFrame} with a three-level column \code{MultiIndex}:
\[
    \Bigl(\underbrace{\text{block}}_{\text{e.g.\ ``chemical''}},\;
          \underbrace{\text{level}}_{\text{``raw''}},\;
          \underbrace{\text{variable}}_{\text{e.g.\ ``Cu''}}\Bigr).
\]
The helper \code{extract\_block(data, block, level)} returns a flat
\code{DataFrame} containing only the columns belonging to the requested
block and level.  Four blocks are used:

\begin{center}
\begin{tabular}{lp{9cm}}
\toprule
Block key & Content \\
\midrule
\code{chemical / raw} & 16 sediment-chemistry concentrations \\
\code{taxa / raw} & 16 benthic invertebrate taxa in \textbf{octave scale} \\
\code{environmental / raw} & Physical habitat variables (depth, DO, temperature,
      MPS, LOI, velocity) \\
\code{sample\_info / raw} & Site identifiers, coordinates, waterbody \\
\bottomrule
\end{tabular}
\end{center}


\subsection{Key Symbols and Variable Types}
\label{sec:data:symbols}

Throughout, we index \(n\) sites by \(i\) and suppress the site subscript
where unambiguous.

\begin{center}
\renewcommand{\arraystretch}{1.25}
\begin{tabularx}{\textwidth}{l l l X}
\toprule
Symbol & Dimensions & Code name & Description \\
\midrule
$\mat{C}$
    & $\R^{n \times 16}$
    & \code{pollution\_raw}
    & Raw sediment-chemistry matrix (sites $\times$ 16 variables) \\
$\tilde{\mat{C}}$
    & $\R^{n \times 16}$
    & \code{pollution\_transformed}
    & $\log_2$-transformed chemistry: $\tilde{c}_{ij}=\log_2(1+c_{ij})$ \\
$\mat{O}$
    & $\R^{n \times p}$
    & \code{taxa\_all}
    & Taxa matrix in \emph{octave} scale ($p=16$ study taxa) \\
$\mat{P}$
    & $\R^{n \times p}$
    & \code{taxa\_relabd}
    & Relative-abundance matrix recovered from octave; rows sum to~1 \\
$\mat{E}$
    & $\R^{n \times q}$
    & \code{env\_ref} or \code{env\_data}
    & Environmental predictor matrix ($q=5$ unless otherwise stated) \\
$\mathbf{s}$
    & $\R^{n}$
    & \code{Pollution\_Score}
    & Composite pollution score (Stage~1 output) \\
$\mathbf{g}$
    & $\{1,\dots,K\}^{n}$
    & \code{Predicted\_Cluster}
    & Habitat-cluster label for each site (Stage~3 output) \\
$\mathbf{x}$
    & $\R^{n \times 2}$
    & \code{coords\_df[NMDS1, NMDS2]}
    & NMDS ordination coordinates (Stage~4) \\
$\mathbf{z}$
    & $[0,1]^{n}$
    & \code{ZCI}
    & Zoobenthic Community Index (Stage~4 output) \\
\bottomrule
\end{tabularx}
\end{center}


\subsection{Characters of the Raw Data}
\label{sec:data:raw}

\subsubsection{Taxa Data in Octave Scale}
\label{sec:data:octave}

The 16 taxa variables are stored in \textbf{octave-transformed} form.
The octave scale encodes the original proportional abundance $p_{ij}$
(taxon $j$ at site~$i$) as
\begin{equation}
    o_{ij} \;=\; \log_{2}\bigl(100\,(p_{ij} + 0.01)\bigr).
    \label{eq:octave_forward}
\end{equation}
Hence $o_{ij}$ increases monotonically with $p_{ij}$ but is not a count;
absolute abundances are \textbf{not recoverable} from the data.

The inverse used by the pipeline (\code{octave\_to\_relative\_abundance()}
in module \code{zci.core.transforms}) is
\begin{equation}
    \hat{p}_{ij} \;=\; \max\!\bigl(0,\; 2^{o_{ij}} - 0.0625\bigr),
    \qquad
    p_{ij}^{\mathrm{RA}} \;=\; \frac{\hat{p}_{ij}}{\sum_{j'}\hat{p}_{ij'}},
    \label{eq:octave_inverse}
\end{equation}
producing a row-normalised relative-abundance matrix $\mat{P}$
($p_{i\cdot}^{\mathrm{RA}}\in[0,1]$, $\sum_j p_{ij}^{\mathrm{RA}} = 1$).

\paragraph{Usage in downstream stages.}
\begin{itemize}[nosep]
    \item \textbf{Stage~2 clustering} uses octave values directly (identity
          transform) when \code{taxa\_transform="octave"} is specified.
    \item \textbf{Stage~4 NMDS / ZCI} converts to relative abundance via
          \eqref{eq:octave_inverse} so that the Bray--Curtis dissimilarity
          operates on compositional proportions.
    \item \textbf{ANOVA} (hindsight step) presents per-taxon means on the
          un-transformed octave scale for interpretability.
    \item \textbf{RDA} can receive either octave or Hellinger-transformed
          taxa (configurable via \code{taxa\_transform}).
\end{itemize}


\subsubsection{Mean Particle Size — MPS(Phi)}
\label{sec:data:mps}

Sediment grain size is stored as \textbf{MPS\,(Phi)}, the mean particle
size after the Krumbein phi transformation:
\begin{equation}
    \phi \;=\; -\log_{2}(\text{particle size in mm}).
    \label{eq:phi}
\end{equation}
The column \code{MPS\,(Phi)} exists in the input workbook as a
pre-computed field; the pipeline reads it directly and does not
recompute $\phi$ from raw particle-size measurements.  It enters
Stage~RDA and Stage~3 (LDA) as one of the environmental predictors
$\mat{E}$.


\subsubsection{Selection of 16 Chemical Variables}
\label{sec:data:chemicals}

The 2004/5 corridor-wide survey measured 19 elements (18 metals plus
arsenic); the 1991 and 1999 Detroit River surveys reported 13 and 19
elements respectively.  The intersection yielded \textbf{13 elements}
(12~metals plus arsenic) that were common to all three surveys.

All three surveys reported PCB congener concentrations, but
methodology and detection limits changed substantially between 1990 and
2004, making congener-by-congener analysis unreliable.  Therefore the
\textbf{sum of PCBs} ($\Sigma$PCBs, recorded as \code{total PCB}) was
used instead.

Pesticide and petrochemical data were variably available across
surveys.  Two single-variable surrogates were chosen:
\begin{itemize}[nosep]
    \item \textbf{p,p'-DDE}: a representative insecticide degradation
          product (surrogate for agricultural pesticides).
    \item \textbf{OCS} (octachlorostyrene): a petrochemical byproduct
          (surrogate for industrial petrochemical contamination).
\end{itemize}

This yields a total of $13 + 1 + 1 + 1 = 16$ chemical variables compiled
into the site-by-contaminant matrix $\mat{C}$.  The exact column names
used in the code (constant \code{POLLUTION\_VARS\_2008} in module
\code{zci.pipeline.pollution\_assessment}) are:

\begin{quote}
\code{Co, Al, Ni, Mn, Fe, Cr, Cu, Hg, Pb, Zn,
      total PCB, Cd, OCS, p,p'-DDE, As, Ca.}
\end{quote}


% ═════════════════════════════════════════════════════════════════════════════
\section{Core Analytical Stream}
\label{sec:core}
% ═════════════════════════════════════════════════════════════════════════════

% ────────────────────────────────────────────────────────────────────────────
\subsection{Stage 1 — Pollution PCA}
\label{sec:core:s1}
% ────────────────────────────────────────────────────────────────────────────

\paragraph{Module.}
\code{zci.pipeline.pollution\_assessment.pollution\_pca\_pipeline()}.

\paragraph{Input.}
Chemical block $\mat{C}\in\R^{n\times 16}$
(\code{extract\_block(data, "chemical", "raw")[POLLUTION\_VARS\_2008]}).

\paragraph{Step 1 — Log transformation.}
Each concentration is log-transformed element-wise
(\code{zci.core.transforms.log2\_transform}):
\begin{equation}
    \tilde{c}_{ij} \;=\; \log_2(1 + c_{ij}),
    \qquad i=1,\dots,n,\;\; j=1,\dots,16.
    \label{eq:log2}
\end{equation}
This yields the transformed matrix
$\tilde{\mat{C}}\in\R^{n\times 16}$
(\code{pollution\_transformed: pd.DataFrame}).

\paragraph{Step 2 — PCA.}
Principal Component Analysis is performed on $\tilde{\mat{C}}$
(\code{zci.core.pca.run\_pca}):
\begin{enumerate}[nosep]
    \item Centre: $\mat{T} = \tilde{\mat{C}} - \mathbf{1}\,\bar{\tilde{\mathbf{c}}}^{\!\top}$.
    \item SVD:  $\mat{T} \approx \mat{U}_k\,\mat{\Sigma}_k\,\mat{V}_k^{\top}$
          with $k$ retained components (default $k=5$).
    \item \textbf{Scaled loadings}:
          $L_{jm} = v_{jm}\sqrt{\lambda_m}$,
          where $\lambda_m$ is the $m$-th eigenvalue and $v_{jm}$ the
          corresponding eigenvector element.
          (\code{result.loadings: pd.DataFrame}, shape $16\times k$.)
    \item \textbf{Site scores}:
          $\mat{S}_{\mathrm{raw}} = \mat{T}\,\mat{V}_k \in\R^{n\times k}$
          (\code{result.scores\_raw}).
    \item \textbf{Standardised scores}: each column of $\mat{S}_{\mathrm{raw}}$
          is min--max rescaled to $[0,1]$
          (\code{result.scores: pd.DataFrame}):
          \begin{equation}
              s_{im}^{*} = \frac{s_{im} - \min_i s_{im}}
                                {\max_i s_{im} - \min_i s_{im}}.
              \label{eq:minmax}
          \end{equation}
\end{enumerate}

\noindent
\textbf{Note.}  The PCA is fit on $\tilde{\mat{C}}$ \emph{without} prior
unit-variance scaling; variables with larger log-scale variance therefore
have more influence on the leading PCs.  This is a deliberate design
choice, not an oversight, but it means that the pollution score is
weighted towards high-variance contaminants.

\paragraph{Step 3 — Composite Pollution Score.}
Selected PCs (by default all $k$) are individually min--max-scaled and
then summed with equal weights
(\code{zci.core.transforms.composite\_pollution\_score}):
\begin{equation}
    s_i \;=\; \sum_{m=1}^{k} w_m \cdot s_{im}^{*},
    \qquad w_m = 1 \;\;\forall\, m \text{ (default)}.
    \label{eq:ps}
\end{equation}
Result: \code{Pollution\_Score: pd.Series} ($\mathbf{s}\in\R^n$).
Higher values indicate greater contamination.

\paragraph{Artifact written.}
\code{01\_updated\_data.xlsx} — the original workbook augmented with
columns \code{PC1}--\code{PC$k$} and \code{Pollution\_Score}, under the
block header
\code{("01\_pollution\_assessment", "raw", *)}.


% ────────────────────────────────────────────────────────────────────────────
\subsection{Stage 2 — Taxa Assemblage Clustering}
\label{sec:core:s2}
% ────────────────────────────────────────────────────────────────────────────

\paragraph{Module.}
\code{zci.pipeline.taxa\_assemblage.taxa\_assemblage\_pipeline()}.

\paragraph{Input.}
\begin{itemize}[nosep]
    \item Taxa block $\mat{O}\in\R^{n\times 16}$ (octave scale).
    \item Pollution Score $\mathbf{s}$ from Stage~1 artifact.
\end{itemize}

\paragraph{Step 1 — Reference-site selection.}
A Boolean mask selects the least-polluted fraction of sites
(\code{zci.core.clustering.select\_reference\_sites}):
\begin{equation}
    \mathcal{R} \;=\; \bigl\{\, i : s_i \;\le\; s_{(n_{\mathrm{ref}})}\,\bigr\},
    \qquad
    n_{\mathrm{ref}} = \lfloor 0.20\,n \rfloor,
    \label{eq:refmask}
\end{equation}
where $s_{(n_{\mathrm{ref}})}$ is the $n_{\mathrm{ref}}$-th smallest
value of~$\mathbf{s}$.  Default: bottom 20\,\% of pollution scores.
Result: \code{ref\_mask: pd.Series[bool]}.

\paragraph{Step 2 — Taxa transform.}
When \code{taxa\_transform="octave"} (the default), the taxa matrix is
used as-is; no numerical transformation is applied.  The function
\code{zci.core.transforms.octave\_transform()} simply returns a copy.

\paragraph{Step 3 — Ward's hierarchical clustering.}
\code{zci.core.clustering.ward\_cluster(taxa\_ref, n\_clusters=$K$)}
performs agglomerative clustering with Ward's linkage (minimum
within-cluster variance) and Euclidean distance:
\begin{equation}
    d_{\mathrm{merge}}(A,B) \;=\;
    \sqrt{\frac{2\,|A|\,|B|}{|A|+|B|}}\,
    \|\bar{\mathbf{o}}_A - \bar{\mathbf{o}}_B\|,
    \label{eq:ward}
\end{equation}
where $\bar{\mathbf{o}}_A$ denotes the centroid of cluster~$A$ in
octave space.  The dendrogram is cut at $K$ clusters (default $K=3$
for the corridor-wide run, later merged to $K=2$ in the hindsight step).

Result: \code{labels\_ref: pd.Series} (1-indexed cluster labels for
reference sites only; non-reference sites are \code{NaN}).

\paragraph{Artifact written.}
\code{02\_updated\_data.xlsx} — block
\code{("02\_taxa\_assemblage", "raw", "Cluster")} populated for
$|\mathcal{R}|$ reference sites.


% ────────────────────────────────────────────────────────────────────────────
\subsection{Hindsight Relabel and ANOVA}
\label{sec:core:hindsight}
% ────────────────────────────────────────────────────────────────────────────

\paragraph{Module.}
\code{run\_hindsight\_relabel.py}; ANOVA core in
\code{zci.core.anova.anova\_table()}.

\paragraph{Cluster relabelling.}
After inspecting Stage~2 dendrograms and cluster membership, smaller
clusters may be merged.  In the corridor-wide run the mapping
$\{3 \to 1\}$ is applied, effectively merging the smallest cluster into
Cluster~1, yielding $K=2$ habitat groups.

\paragraph{One-way ANOVA.}
For each environmental variable $e$ (from $\mat{E}$) and each taxon
(from $\mat{O}$, restricted to reference sites), a one-way ANOVA across
the $K$ clusters is computed:
\begin{equation}
    F = \frac{\mathrm{SS}_{\text{between}} / (K-1)}
             {\mathrm{SS}_{\text{within}} / (n_{\mathrm{ref}}-K)},
    \label{eq:anova}
\end{equation}
where
\[
    \mathrm{SS}_{\text{between}} = \sum_{g=1}^{K} n_g\,
        (\bar{e}_g - \bar{e})^2,
    \qquad
    \mathrm{SS}_{\text{within}} = \sum_{g=1}^{K}
        \sum_{i\in g} (e_i - \bar{e}_g)^2.
\]
The $p$-value is obtained from \code{scipy.stats.f\_oneway}.

The code supports an optional pre-ANOVA transform controlled by the
parameter \code{transform}:
\begin{itemize}[nosep]
    \item \code{"none"} — no transformation (default for the corridor run).
    \item \code{"log"} — $\ln(1+x)$ applied to each variable.
    \item \code{"boxcox"} — Box--Cox power transformation (with fallback
          to $\ln(1+x)$ on failure).
\end{itemize}
Per-cluster means $\pm$ SEM are reported on the \emph{untransformed}
scale regardless of the ANOVA transform used.

\paragraph{Outputs.}
\begin{itemize}[nosep]
    \item \code{anova\_env.xlsx}, \code{anova\_taxa.xlsx} — ANOVA summary
          tables with $F$, $p$, cluster means $\pm$ SEM.
    \item \code{cluster\_panel.png} — map + environmental bars + taxa bars.
    \item \code{02\_hindsight\_updated\_data.xlsx} — relabelled artifact
          (block column \code{Cluster} updated).
\end{itemize}


% ────────────────────────────────────────────────────────────────────────────
\subsection{Stage 3 — LDA Habitat Classification}
\label{sec:core:s3}
% ────────────────────────────────────────────────────────────────────────────

\paragraph{Module.}
\code{zci.pipeline.lda\_classification.lda\_pipeline()};
core in \code{zci.core.lda}.

\paragraph{Input.}
\begin{itemize}[nosep]
    \item Environmental matrix $\mat{E}_{\mathcal{R}}$ for reference sites,
          comprising the $q=5$ variables:
          \code{Measured Depth (m)}, \code{Water DO Bottom (mg/L)},
          \code{Temperature (oC)}, \code{MPS (Phi)}, \code{LOI (\%)}.
    \item Cluster labels $\mathbf{g}_{\mathcal{R}}$ from the hindsight
          artifact.
    \item Pollution Score $\mathbf{s}$ (for reference-site identification).
\end{itemize}

\paragraph{Step 1 — Standardisation.}
When \code{standardize=True} (default), each environmental variable is
$z$-scored using \code{sklearn.preprocessing.StandardScaler}:
\begin{equation}
    e_{ij}^{*} = \frac{e_{ij} - \bar{e}_j}{\sigma_j},
    \label{eq:zscore}
\end{equation}
fitted \emph{on the reference sites only}.  The same scaler object is
stored in the \code{LDAFit} container and reapplied to non-reference
sites at prediction time
(\code{zci.core.lda.predict\_sites}).

\paragraph{Step 2 — LDA fitting.}
\code{sklearn.discriminant\_analysis.LinearDiscriminantAnalysis} is fit on
$(\mat{E}_{\mathcal{R}}^{*},\,\mathbf{g}_{\mathcal{R}})$.  For $K$
classes and $q$ predictors the model extracts $\min(K-1,\,q)$ linear
discriminant axes, each a linear combination of the standardised
environmental variables that maximises the between-class to
within-class variance ratio:
\begin{equation}
    \mathrm{LD}_m = \mathbf{w}_m^{\!\top}\,\mathbf{e}^{*},
    \qquad
    \mathbf{w}_m = \arg\max_{\mathbf{w}} \;
    \frac{\mathbf{w}^{\!\top}\mat{S}_B\,\mathbf{w}}
         {\mathbf{w}^{\!\top}\mat{S}_W\,\mathbf{w}},
    \label{eq:lda}
\end{equation}
where $\mat{S}_B$ and $\mat{S}_W$ are the between- and within-class
scatter matrices respectively.

\paragraph{Step 3 — Variable importance (Wilks' Lambda).}
\code{zci.core.lda.wilks\_lambda\_importance()} performs a drop-one
analysis.  For each variable $j$:
\begin{align}
    \Lambda_{\text{full}} &= \frac{|\mat{W}|}{|\mat{T}|},
    \label{eq:wilks_full} \\
    \Lambda_{-j} &= \text{Wilks' }\Lambda\text{ with variable }j\text{ removed},
    \label{eq:wilks_drop} \\
    \Delta\Lambda_j &= \Lambda_{-j} - \Lambda_{\text{full}},
    \label{eq:wilks_delta}
\end{align}
and an approximate $F$-statistic is computed:
\begin{equation}
    F_j = \frac{\Delta\Lambda_j / \Lambda_{\text{full}}}
               {(K-1)^{-1}\,(n_{\mathrm{ref}} - K - q + 1)^{-1}}.
    \label{eq:wilks_f}
\end{equation}
\noindent
Overall model significance is assessed via Bartlett's $\chi^2$
approximation:
\begin{equation}
    \chi^2 = -\Bigl(n_{\mathrm{ref}} - 1
             - \frac{q+K}{2}\Bigr)\,\ln\Lambda_{\text{full}},
    \qquad
    \text{df} = q\,(K-1).
    \label{eq:bartlett}
\end{equation}

\paragraph{Step 4 — Monte Carlo cross-validation (MCCV).}
\code{zci.core.lda.monte\_carlo\_cv()} repeats stratified 80/20
train/test splits for $N_{\mathrm{cv}}$ iterations (default 1\,000):
\begin{enumerate}[nosep]
    \item \emph{Globally} standardise (z-score) the reference environmental
          data.
    \item For each iteration: draw a stratified 80\,\% training set,
          fit LDA, predict on the held-out 20\,\%, record accuracy.
    \item Report mean $\pm$ std accuracy, per-class precision / recall /
          F1, and an aggregated confusion matrix.
\end{enumerate}

\noindent
\textbf{Caveat}: the z-score standardisation in the current code is
applied \emph{globally before} the CV splits, not refitted inside each
fold.  This is a minor data-leakage issue; reported CV accuracies may be
slightly optimistic.

\paragraph{Step 5 — Classification of non-reference sites.}
\code{zci.core.lda.predict\_sites()} applies the fitted LDA (with the
stored scaler) to the environmental data of all non-reference sites,
producing:
\begin{itemize}[nosep]
    \item \code{Predicted\_Cluster: pd.Series} — the habitat label $g_i$
          for each non-reference site.
    \item \code{probabilities: pd.DataFrame} — posterior class
          probabilities.
    \item An \code{Is\_Reference: pd.Series[bool]} flag.
\end{itemize}

\paragraph{Artifact written.}
\code{03\_updated\_data.xlsx} — block
\code{("03\_lda\_classification", "raw", *)}, adding columns
\code{Predicted\_Cluster}, \code{Is\_Reference}, per-class probabilities,
and LD coordinates.


% ────────────────────────────────────────────────────────────────────────────
\subsection{Stage 4 — Bray--Curtis NMDS and ZCI Construction}
\label{sec:core:s4}
% ────────────────────────────────────────────────────────────────────────────

\paragraph{Module.}
\code{zci.pipeline.bray\_curtis\_nmds.nmds\_pipeline()};
core in \code{zci.core.nmds} and \code{zci.core.zci\_scores}.

\paragraph{Input.}
\begin{itemize}[nosep]
    \item Taxa block $\mat{O}$ (octave scale) from the original workbook.
    \item Pollution Score $\mathbf{s}$ from Stage~1 artifact.
    \item Predicted Cluster $\mathbf{g}$ and \code{Is\_Reference} from
          Stage~3 artifact.
\end{itemize}

\paragraph{Step 1 — Octave to relative abundance.}
The taxa matrix is converted to relative abundance via
\code{octave\_to\_relative\_abundance()} as given in
\eqref{eq:octave_inverse}, producing $\mat{P}\in\R^{n\times p}$
(\code{taxa\_relabd: pd.DataFrame}).

\paragraph{Step 2 — Endpoint construction (per cluster).}
For each cluster $g$, the $n_{\mathrm{ep}}$ sites with the lowest
Pollution Score form the \textbf{reference endpoint} and the
$n_{\mathrm{ep}}$ highest form the \textbf{degraded endpoint}
(\code{zci.core.nmds.build\_endpoints}):
\begin{equation}
    \bar{\mathbf{p}}_{\text{ref}} =
        \frac{1}{n_{\mathrm{ep}}} \sum_{i\in\mathcal{E}_{\text{ref}}}
        \mathbf{p}_i^{\mathrm{RA}},
    \qquad
    \bar{\mathbf{p}}_{\text{deg}} =
        \frac{1}{n_{\mathrm{ep}}} \sum_{i\in\mathcal{E}_{\text{deg}}}
        \mathbf{p}_i^{\mathrm{RA}},
    \label{eq:endpoints}
\end{equation}
where $\mathcal{E}_{\text{ref}}$ and $\mathcal{E}_{\text{deg}}$ are the
index sets of the $n_{\mathrm{ep}}$ extreme sites.  These two
``synthetic sites'' are appended to the cluster's taxa matrix.

\paragraph{Step 3 — Bray--Curtis dissimilarity.}
The pairwise Bray--Curtis dissimilarity between every pair of sites
(including endpoints) in the augmented relative-abundance matrix is
computed (\code{scipy.spatial.distance.pdist(..., metric="braycurtis")}):
\begin{equation}
    d_{\BC}(i,i') \;=\;
    \frac{\displaystyle\sum_{j=1}^{p} |p_{ij}^{\mathrm{RA}} - p_{i'j}^{\mathrm{RA}}|}
         {\displaystyle\sum_{j=1}^{p} (p_{ij}^{\mathrm{RA}} + p_{i'j}^{\mathrm{RA}})},
    \label{eq:braycurtis}
\end{equation}
yielding $\mat{D}\in[0,1]^{(n_g+2)\times(n_g+2)}$, where $n_g$ is the
number of real sites in cluster~$g$.

\emph{This is the first place where Bray--Curtis is computed; it feeds
the NMDS ordination.}

\paragraph{Step 4 — Iterative non-metric MDS.}
\code{zci.core.nmds.iterative\_nmds()} embeds $\mat{D}$ into
$\R^2$ via non-metric multidimensional scaling (NMDS):
\begin{enumerate}[nosep]
    \item \textbf{Pass 1}: \code{sklearn.manifold.MDS} with
          \code{metric=False} (non-metric) and \code{n\_init=10}
          random starts.
    \item \textbf{Passes 2--3}: re-initialise MDS with the previous
          solution; \code{n\_init=4} additional random starts per
          refinement.
    \item Total: 3~iterations of SMACOF (Scaling by MAjorising a
          Complicated Function).
\end{enumerate}
The algorithm minimises the \emph{stress} criterion:
\begin{equation}
    \mathrm{Stress} = \sqrt{
        \frac{\sum_{i<i'}\bigl(f(d_{\BC}(i,i')) -
              \|\mathbf{x}_i - \mathbf{x}_{i'}\|\bigr)^2}
             {\sum_{i<i'} \|\mathbf{x}_i - \mathbf{x}_{i'}\|^2}
    },
    \label{eq:stress}
\end{equation}
where $f$ is a monotone regression function (the ``non-metric'' part).

\paragraph{Step 5 — PCA rotation and axis flip.}
\code{zci.core.nmds.pca\_rotate()} rotates the NMDS coordinates so
NMDS1 captures the maximum variance.  Then
\code{zci.core.nmds.flip\_axis()} checks the Pearson correlation $r$
between NMDS1 and $\mathbf{s}$ (Pollution Score) for the real sites:
\begin{equation}
    \text{if } r(\mathrm{NMDS1},\,\mathbf{s}) > 0
    \;\;\Longrightarrow\;\;
    \mathrm{NMDS1} \leftarrow -\mathrm{NMDS1}.
    \label{eq:flip}
\end{equation}
This enforces a convention where NMDS1 is \emph{negatively} correlated
with pollution (cleaner sites on the positive end).

\textbf{Interpretive caveat}: this axis-flip is enforced by design, so
the resulting negative NMDS1--Pollution correlation is partly an artefact
of the orientation convention and should not be interpreted as a
discovered ecological relationship.

\paragraph{Step 6 — Weighted-average species scores.}
\code{zci.core.nmds.weighted\_average\_scores()} assigns each taxon a
position in NMDS space:
\begin{equation}
    \mathrm{WA}_{j} = \frac{\sum_{i=1}^{n_g} p_{ij}^{\mathrm{RA}}\,\mathbf{x}_i}
                            {\sum_{i=1}^{n_g} p_{ij}^{\mathrm{RA}}}.
    \label{eq:wa}
\end{equation}

\paragraph{Step 7 — ZCI computation.}

The pipeline currently uses the \textbf{BC-Direct} method
(\code{zci.core.zci\_scores.zci\_bc\_direct}) as the default ZCI:
\begin{equation}
    z_i \;=\; \frac{d_{\BC}(i,\,\mathcal{E}_{\text{deg}})}
                    {d_{\BC}(i,\,\mathcal{E}_{\text{ref}})
                     + d_{\BC}(i,\,\mathcal{E}_{\text{deg}})},
    \label{eq:zci_bcdirect}
\end{equation}
where $d_{\BC}(i,\,\mathcal{E})$ is the Bray--Curtis distance from site~$i$'s
relative-abundance vector to the centroid of the endpoint set $\mathcal{E}$
(computed via \code{scipy.spatial.distance.cdist}).

\emph{This is the second place where Bray--Curtis is applied; it operates
directly on relative-abundance vectors and does \textbf{not} use the
NMDS coordinates.}  Values near~1 indicate a community similar to the
reference endpoint; values near~0 indicate a degraded community.

The default per-cluster endpoint sizes in the corridor run are:
\begin{center}
\begin{tabular}{ccc}
\toprule
Cluster & $n_{\mathrm{ep}}$ (ZCI) & Method \\
\midrule
1 & 15 & BC-Direct \\
2 & 3  & BC-Direct \\
\bottomrule
\end{tabular}
\end{center}

ZCI--Pollution correlations (Pearson and Spearman) are computed via
\code{zci.core.zci\_scores.zci\_correlation()} and saved in
\code{zci\_summary.xlsx}.

\paragraph{Alternative ZCI methods (available but not used by default).}
The module \code{zci.core.zci\_scores} also implements:
\begin{itemize}[nosep]
    \item \textbf{Distance}: $z_i = d_{\text{deg}} / (d_{\text{ref}} + d_{\text{deg}})$
          using Euclidean distances in NMDS space.
    \item \textbf{Projection}: project sites onto the directed
          DEG$\to$REF line in NMDS space, normalised to $[0,1]$.
    \item \textbf{Centroid-Proj}: project onto a centroid-defined axis
          using real-site pollution extremes (rather than synthetic endpoints).
\end{itemize}

\paragraph{Artifact written.}
\code{04\_updated\_data.xlsx} — block
\code{("04\_bray\_curtis\_nmds", "raw", *)}, adding columns
\code{NMDS1}, \code{NMDS2}, \code{ZCI}, \code{ZCI\_Method}.


% ────────────────────────────────────────────────────────────────────────────
\subsection{Stage 5 — Piecewise Quantile Regression}
\label{sec:core:s5}
% ────────────────────────────────────────────────────────────────────────────

\paragraph{Module.}
\code{zci.pipeline.piecewise\_qr\_pipeline.pqr\_pipeline()};
core in \code{zci.core.piecewise\_qr}.

\paragraph{Input.}
\begin{itemize}[nosep]
    \item $\mathbf{s}$ (Pollution Score) from Stage~1.
    \item $\mathbf{g}$ (Predicted Cluster) from Stage~3.
    \item $\mathbf{z}$ (ZCI) from Stage~4.
\end{itemize}

\paragraph{Model.}
For each cluster and a grid of quantile levels
$\tau\in\{0.10,0.15,\dots,0.90\}$, a piecewise linear quantile
regression is fit:
\begin{equation}
    Q_\tau(\ZCI \mid \PS = x)
    \;=\; \beta_0 + \beta_1\,x + \beta_2\,(x - \psi)_+
    \label{eq:pqr}
\end{equation}
where $(x-\psi)_+ = \max(0,\, x-\psi)$ is a hinge function at
breakpoint~$\psi$.

\paragraph{Estimation.}
The parameters $(\beta_0,\beta_1,\beta_2,\psi)$ are estimated
jointly by a \textbf{profile grid search}
(\code{zci.core.piecewise\_qr.\_profile\_search}):
\begin{enumerate}[nosep]
    \item Fix a candidate breakpoint $\psi_c$ on a grid spanning the
          10th--90th percentile of~$\mathbf{s}$.
    \item For each $\psi_c$, build the design matrix
          $\mat{X}(\psi_c) = [\mathbf{1},\;\mathbf{s},\;(\mathbf{s}-\psi_c)_+]$
          and solve the linear quantile regression sub-problem by
          Linear Programming (LP via \code{scipy.optimize.linprog}):
          \begin{equation}
              \min_{\boldsymbol{\beta},\mathbf{u},\mathbf{v}}\;
              \tau\,\mathbf{1}^{\!\top}\mathbf{u}
              + (1-\tau)\,\mathbf{1}^{\!\top}\mathbf{v}
              \quad\text{s.t.}\quad
              \mat{X}\boldsymbol{\beta} + \mathbf{u} - \mathbf{v} = \mathbf{z},
              \;\; \mathbf{u},\mathbf{v}\ge\mathbf{0}.
              \label{eq:qr_lp}
          \end{equation}
    \item Select the $\psi_c$ with the lowest check loss
          $\rho_\tau(\mathbf{r}) = \sum_i r_i\,[\tau - 1(r_i < 0)]$.
\end{enumerate}

\paragraph{Wild-bootstrap confidence intervals.}
\code{zci.core.piecewise\_qr.wild\_bootstrap\_ci()} follows
Feng, He \& Hu (2011):
\begin{enumerate}[nosep]
    \item Fit the original model; obtain residuals $\hat{r}_i = z_i - \hat{Q}_\tau(s_i)$.
    \item For $b = 1,\dots,B$ (default $B=500$):
    \begin{enumerate}[nosep]
        \item Draw Mammen's two-point weights $w_i$:
        \[
            w_i = \begin{cases}
                -\frac{\sqrt{5}-1}{2} & \text{w.p. } \frac{\sqrt{5}+1}{2\sqrt{5}}, \\[4pt]
                \phantom{-}\frac{\sqrt{5}+1}{2} & \text{w.p. } \frac{\sqrt{5}-1}{2\sqrt{5}}.
            \end{cases}
        \]
        \item Construct perturbed response:
              $z_i^{*} = \hat{Q}_\tau(s_i) + w_i\,\hat{r}_i$.
        \item Re-fit the piecewise QR on $(s_i,\,z_i^{*})$ to obtain
              $\hat{\boldsymbol{\beta}}^{(b)},\,\hat{\psi}^{(b)}$.
    \end{enumerate}
    \item Confidence intervals: $\alpha/2$ and $1-\alpha/2$ quantiles
          of the bootstrap distribution of each parameter
          (default: 90\,\% CI).
\end{enumerate}

\paragraph{Note on current status.}
In the corridor-wide runner
(\code{run\_SCDRS\_full\_pipeline.py}), Stage~5 is currently
\textbf{commented out} (lines are present but disabled); it is active
in the Detroit River runner
(\code{run\_DR\_full\_pipeline.py}).  The mathematical description above
applies to both.

\paragraph{Artifact written.}
\code{05\_updated\_data.xlsx} (when the stage is enabled).


% ═════════════════════════════════════════════════════════════════════════════
\section{Complementary / Validation Methods}
\label{sec:complementary}
% ═════════════════════════════════════════════════════════════════════════════

% ────────────────────────────────────────────────────────────────────────────
\subsection{Stage RDA — Redundancy Analysis}
\label{sec:complementary:rda}
% ────────────────────────────────────────────────────────────────────────────

\paragraph{Module.}
\code{zci.pipeline.rda\_analysis.rda\_pipeline()};
core in \code{zci.core.rda.RDA}.

\paragraph{Purpose.}
RDA is a \textbf{complementary validation} method, not part of the main
ZCI construction stream.  It quantifies how much of the taxa variation
among reference sites is linearly explained by the environmental
predictors, and it identifies which environmental variables are most
important.

\paragraph{Input.}
\begin{itemize}[nosep]
    \item Taxa matrix $\mat{Y}$ for reference sites ($\mat{O}_{\mathcal{R}}$
          in octave scale, or Hellinger-transformed).
    \item Environmental matrix $\mat{X}$ ($q=5$ variables, optionally
          $\ln(1+x)$-transformed and/or $z$-scored).
    \item Reference mask from Stage~1 Pollution Score.
\end{itemize}

\paragraph{Step 1 — Multivariate OLS.}
Centre both matrices; fit OLS:
\begin{equation}
    \hat{\mat{Y}} = \mat{X}_c\,\hat{\mat{B}},
    \qquad
    \hat{\mat{B}} = (\mat{X}_c^{\!\top}\mat{X}_c)^{-1}\mat{X}_c^{\!\top}\mat{Y}_c.
    \label{eq:rda_ols}
\end{equation}

\paragraph{Step 2 — PCA on fitted values.}
Eigendecompose the covariance of $\hat{\mat{Y}}$:
\begin{equation}
    \frac{1}{n-1}\,\hat{\mat{Y}}^{\!\top}\hat{\mat{Y}}
    = \mat{A}\,\mat{\Lambda}\,\mat{A}^{\!\top},
    \label{eq:rda_pca}
\end{equation}
yielding constrained eigenvalues $\lambda_1 \ge \lambda_2 \ge \cdots$
and eigenvectors (species scores).

\paragraph{Step 3 — Scores.}
\begin{itemize}[nosep]
    \item Site scores (scaling~1):
          $\mat{U} = \hat{\mat{Y}}\,\mat{A}\,\mat{\Lambda}^{-1/2}$.
    \item Species scores:
          $\mat{V} = \mat{A}\,\mat{\Lambda}^{1/2}$.
    \item Biplot scores (correlation of $\mat{X}_c$ with $\mat{U}$):
          used for environmental-arrow overlays.
\end{itemize}

\paragraph{Step 4 — Permutation tests.}
Three types of permutation test are implemented, all based on
row-permutation of $\mat{Y}_c$ (or residual permutation for term tests):
\begin{center}
\begin{tabular}{lp{8.5cm}}
\toprule
Test & Null hypothesis \\
\midrule
Global & No linear relationship between $\mat{Y}$ and $\mat{X}$.
         Pseudo-$F = (\text{Inertia}_{\text{con}}/\text{df}_{\text{model}})
         / (\text{Inertia}_{\text{res}}/\text{df}_{\text{res}})$. \\
Per-axis & Each RDA axis individually non-significant.
           $F_m = \lambda_m / (\text{Inertia}_{\text{res}}/\text{df}_{\text{res}})$. \\
Per-term (Freedman--Lane) & Each environmental variable individually
           non-significant, controlling for the others. \\
\bottomrule
\end{tabular}
\end{center}
Default: 999 permutations per test type.

\paragraph{Outputs.}
\code{rda\_axes\_summary.xlsx}, \code{rda\_terms\_summary.xlsx},
\code{rda\_triplot.png}.
No augmented artifact is written — RDA does \emph{not} feed downstream
stages.


% ────────────────────────────────────────────────────────────────────────────
\subsection{Role of ANOVA in the Framework}
\label{sec:complementary:anova}
% ────────────────────────────────────────────────────────────────────────────

ANOVA (\cref{sec:core:hindsight}) is \textbf{supportive}: it validates
that the Ward-clustering-derived habitat groups are ecologically
distinct in terms of environmental conditions and taxa composition.
It does not feed into the quantitative ZCI construction.


% ═════════════════════════════════════════════════════════════════════════════
\section{Artifact and Intermediate-Value Map}
\label{sec:artifact_map}
% ═════════════════════════════════════════════════════════════════════════════

\Cref{tab:artifacts} lists every stage, its input artifact(s), the new
columns added, and the output artifact path.

\begin{landscape}
\begin{longtable}{lp{4.5cm}p{5.5cm}p{5cm}}
\caption{Stage-by-stage artifact chain.  ``Block'' column keys follow the
\code{(block, level, variable)} MultiIndex convention.}
\label{tab:artifacts} \\
\toprule
Stage & Input Artifact(s) & New Columns Added & Output Artifact \\
\midrule
\endfirsthead
\toprule
Stage & Input Artifact(s) & New Columns Added & Output Artifact \\
\midrule
\endhead
\bottomrule
\endfoot

Stage~1 (PCA)
    & Original workbook (chemical/raw)
    & \code{PC1}--\code{PC5}, \code{Pollution\_Score},
      \code{PC1\_sel}--\code{PC5\_sel} (selected PCs)
    & \code{01\_updated\_data.xlsx}
\\[4pt]

Stage~RDA
    & Original workbook + Stage~1 artifact
    & \emph{None} (tables \& figures only)
    & \emph{No artifact; outputs in \code{RDA\_analysis/}}
\\[4pt]

Stage~2 (Clustering)
    & Original workbook + Stage~1 artifact
    & \code{Cluster} (reference sites only),
      \code{Is\_Reference}
    & \code{02\_updated\_data.xlsx}
\\[4pt]

Hindsight
    & Stage~2 artifact + original workbook
    & Relabelled \code{Cluster}
    & \code{02\_hindsight\_updated\_data.xlsx}
\\[4pt]

Stage~3 (LDA)
    & Original workbook + Stage~1 + Hindsight artifacts
    & \code{Predicted\_Cluster}, \code{Is\_Reference},
      per-class probabilities,
      \code{LD1}--\code{LD$(K{-}1)$}
    & \code{03\_updated\_data.xlsx}
\\[4pt]

Stage~4 (NMDS+ZCI)
    & Original workbook + Stage~1 + Stage~3 artifacts
    & \code{NMDS1}, \code{NMDS2}, \code{ZCI}, \code{ZCI\_Method}
    & \code{04\_updated\_data.xlsx}
\\[4pt]

Stage~5 (PQR)
    & Stage~1 + Stage~3 + Stage~4 artifacts
    & \emph{(model parameters; no per-site columns)}
    & \code{05\_updated\_data.xlsx}
\\
\end{longtable}
\end{landscape}


% ═════════════════════════════════════════════════════════════════════════════
\section{Quality Assessment and Potential Issues}
\label{sec:qa}
% ═════════════════════════════════════════════════════════════════════════════

% ────────────────────────────────────────────────────────────────────────────
\subsection{ANOVA Assumption Checking}
\label{sec:qa:anova}
% ────────────────────────────────────────────────────────────────────────────

The one-way ANOVA (\cref{sec:core:hindsight}) assumes:
\begin{enumerate}[nosep]
    \item \textbf{Independence} of observations.
    \item \textbf{Normality} of the residuals within each group.
    \item \textbf{Homogeneity of variances} (homoscedasticity) across groups.
\end{enumerate}

\paragraph{Current status.}
The code in \code{zci.core.anova.anova\_table()} calls
\code{scipy.stats.f\_oneway} after an optional pre-transform
(\code{"none"}, \code{"log"}, or \code{"boxcox"}), but
\textbf{does not} perform any assumption diagnostics:
\begin{itemize}[nosep]
    \item No Shapiro--Wilk or Anderson--Darling test for within-group
          normality.
    \item No Levene or Brown--Forsythe test for equal variances.
    \item No Q--Q plot or residual histogram is generated.
\end{itemize}

\paragraph{Recommended remediation.}
\begin{enumerate}
    \item For each variable $\times$ cluster combination, compute:
    \begin{itemize}[nosep]
        \item Shapiro--Wilk $W$ and $p$-value for within-group normality.
        \item Levene's test (median-based variant, i.e.\ Brown--Forsythe)
              for homogeneity of variances.
    \end{itemize}
    \item When normality or equal-variance assumptions are violated:
    \begin{itemize}[nosep]
        \item Apply a variance-stabilising transform (log, Box--Cox) and
              re-run; the code already supports \code{transform="log"} and
              \code{transform="boxcox"}.
        \item Report Welch's ANOVA (\code{scipy.stats.alexandergovern}
              or the \code{pingouin} library) as a robust alternative.
        \item As a non-parametric backup, report Kruskal--Wallis $H$
              with Dunn's post-hoc tests.
    \end{itemize}
    \item In the methods section, explicitly state which transform
          (if any) was applied, and which variables violated assumptions.
    \item Report effect sizes (eta-squared $\eta^2$) alongside $F$ and $p$
          to assess practical significance.
\end{enumerate}

\paragraph{Impact on results.}
ANOVA results are used to characterise habitat clusters, not to drive the
ZCI computation; thus any bias in $p$-values affects interpretation
tables but not the index itself.  Nevertheless, claiming significance for
specific variables without verifying assumptions weakens the
ecological narrative.


% ────────────────────────────────────────────────────────────────────────────
\subsection{Bray--Curtis and NMDS: Where Each Operates}
\label{sec:qa:bcnmds}
% ────────────────────────────────────────────────────────────────────────────

\paragraph{Clarification of roles.}
Bray--Curtis dissimilarity is applied at \textbf{two distinct points}
in Stage~4:

\begin{enumerate}
    \item \textbf{For NMDS ordination} (\cref{sec:core:s4}, Step~3):
          a full pairwise $\mat{D}$ matrix is computed among all sites
          within a cluster (plus two synthetic endpoints).  This
          $\mat{D}$ is the input to the iterative non-metric MDS
          algorithm, which embeds sites in 2-D ordination space.

    \item \textbf{For BC-Direct ZCI} (\cref{sec:core:s4}, Step~7):
          each site's relative-abundance vector is compared to the
          reference and degraded endpoint \emph{centroids} using
          Bray--Curtis.  This computation is independent of NMDS;
          it operates on the raw $\mat{P}$ matrix directly.
\end{enumerate}

\paragraph{Does the framework correctly reflect contamination-driven
community change?}
\begin{itemize}
    \item The NMDS ordination faithfully represents the community
          \emph{structure} (rank-order of dissimilarities), and its
          stress value (reported per cluster in
          \code{nmds\_summary.xlsx}) quantifies the quality of the
          2-D approximation.
    \item The NMDS1 axis is then \emph{artificially oriented} to
          correlate negatively with pollution (\eqref{eq:flip}).  This
          is a visualisation convenience but means the
          NMDS1--Pollution correlation should not be cited as evidence
          of a contamination gradient.
    \item The \textbf{ZCI itself} (BC-Direct) is computed without any
          reference to the NMDS coordinates.  It directly measures how
          similar each site's community is to the reference vs.\
          degraded endpoint in Bray--Curtis space.  This is a more
          defensible measure of contamination-induced community shift
          because it avoids the information loss inherent in
          projecting to 2-D.
    \item The ZCI--Pollution Score correlation (Pearson and Spearman)
          is computed and reported; a high positive correlation
          supports the interpretation that the index captures
          contamination-driven community change.
\end{itemize}

\paragraph{Potential issues.}
\begin{itemize}[nosep]
    \item \textbf{Endpoint sensitivity}: the $n_{\mathrm{ep}}$ parameter
          controls how many extreme sites define the reference/degraded
          centroids.  Small $n_{\mathrm{ep}}$ (e.g.\ 3 for Cluster~2)
          makes the centroid sensitive to individual sites.
    \item \textbf{NMDS stress}: if stress exceeds $\sim$0.20 the 2-D
          representation is unreliable; report and evaluate per cluster.
    \item \textbf{Cluster-specific ZCI}: the ZCI is computed separately
          within each habitat cluster, so cross-cluster comparisons
          require caution.
\end{itemize}


% ────────────────────────────────────────────────────────────────────────────
\subsection{LDA Cross-Validation Leakage}
\label{sec:qa:lda_leak}
% ────────────────────────────────────────────────────────────────────────────

As noted in \cref{sec:core:s3}, the Monte Carlo CV in
\code{zci.core.lda.monte\_carlo\_cv()} applies
\code{StandardScaler.fit\_transform} to the \emph{entire} reference-site
environmental matrix before splitting into CV folds.
This means every training fold has indirect information about the
test-fold means and variances.  The resulting CV accuracy may be
\textbf{slightly optimistic}.

\paragraph{Recommended fix.}
Move the standardisation inside each CV split: fit the scaler on the
training fold only, then transform the test fold using the training
scaler.


% ═════════════════════════════════════════════════════════════════════════════
\section{TITAN Embedding Plan}
\label{sec:titan}
% ═════════════════════════════════════════════════════════════════════════════

\subsection{Background}
\label{sec:titan:background}

\textbf{TITAN} (Threshold Indicator Taxa ANalysis; Baker \& King, 2010)
detects taxon-specific change points along a continuous environmental
gradient.  It complements the current pipeline by answering a question
the existing stages do not:

\begin{quote}
\emph{At what pollution level does each individual taxon show the
sharpest decline (or increase), and what is the aggregate
community-level threshold?}
\end{quote}

Existing notebook: \code{notebooks/10\_TITAN.ipynb} contains a
simulated-data demonstration and a real-data application (Part~2)
using the corridor-wide pollution-score gradient.

\subsection{TITAN Algorithm Summary}
\label{sec:titan:algorithm}

Given an ordered gradient $x_1 \le x_2 \le \cdots \le x_n$ (Pollution
Score) and a taxa-by-site abundance matrix $\mat{Y}$:

\begin{enumerate}
    \item \textbf{Candidate partitions.}
    For each candidate split value $x_c$ (excluding the 5 smallest and
    5 largest values to avoid edge effects), divide sites into two
    groups: $\{i: x_i \le x_c\}$ and $\{i: x_i > x_c\}$.

    \item \textbf{Indicator Value.}
    For each taxon $j$ at each split~$x_c$, compute the Indicator Value
    (IndVal; Dufr\^{e}ne \& Legendre, 1997):
    \begin{equation}
        \mathrm{IndVal}_j = A_j \times B_j \times 100,
        \label{eq:indval}
    \end{equation}
    where $A_j$ (specificity) is the mean abundance in the target group
    relative to the overall mean, and $B_j$ (fidelity) is the proportion
    of sites in the target group where taxon~$j$ occurs.

    \item \textbf{$z$-score.}
    Compare $\mathrm{IndVal}_j(x_c)$ to a null distribution obtained by
    permuting site labels (250--1000 permutations), yielding
    $z_j(x_c)$.

    \item \textbf{Change-point selection.}
    The change point for taxon $j$ is
    $x_c^{*(j)} = \arg\max_{x_c} z_j(x_c)$.

    \item \textbf{Direction assignment.}
    \begin{itemize}[nosep]
        \item $z^{-}$ (sensitive / clean-water indicator): taxon
              abundance peaks in the low-$x$ group.
        \item $z^{+}$ (tolerant / pollution indicator): taxon abundance
              peaks in the high-$x$ group.
    \end{itemize}

    \item \textbf{Bootstrap validation} ($B_{\mathrm{boot}}$ resamples,
    e.g.\ 500):
    \begin{itemize}[nosep]
        \item \textbf{Purity} $\ge$ 0.95: $\ge$95\,\% of bootstrap
              replicates assign the taxon to the same direction.
        \item \textbf{Reliability} $\ge$ 0.95: $\ge$95\,\% yield a
              significant response ($p \le 0.05$).
    \end{itemize}
    Only taxa passing both filters are retained.

    \item \textbf{Community-level thresholds.}
    Sum the filtered $z$-scores across all reliable taxa in each direction:
    \begin{equation}
        \mathrm{sum}(z^{-})(x) = \sum_{j\in\text{declining}} z_j^{-}(x),
        \qquad
        \mathrm{sum}(z^{+})(x) = \sum_{j\in\text{increasing}} z_j^{+}(x).
        \label{eq:sumz}
    \end{equation}
    The $x$-value where each sum peaks defines the community-level
    change point.
\end{enumerate}

\subsection{Integration Points in the Current Framework}
\label{sec:titan:integration}

TITAN should be run \textbf{per habitat cluster} (using Stage~3
\code{Predicted\_Cluster} labels) with the following inputs:

\begin{center}
\begin{tabular}{lll}
\toprule
TITAN input & Source & Code path \\
\midrule
Gradient $x$ & \code{Pollution\_Score} & Stage~1 artifact \\
Taxa matrix $\mat{Y}$ & Relative-abundance $\mat{P}$
    & \code{octave\_to\_relative\_abundance(taxa\_all)} \\
Cluster grouping & \code{Predicted\_Cluster} & Stage~3 artifact \\
\bottomrule
\end{tabular}
\end{center}

\paragraph{Execution.}
TITAN2 is an R package; the existing notebook
(\code{notebooks/10\_TITAN.ipynb}) calls it via \code{rpy2} magic cells.
A production integration would:
\begin{enumerate}[nosep]
    \item Load Stage~1, Stage~3, and Stage~4 artifacts plus the original
          taxa block.
    \item For each cluster, filter sites and pass the Pollution Score
          vector and taxa matrix to \code{TITAN2::titan()} in R.
    \item Retrieve \code{sppmax} (taxon-level change points, purity,
          reliability) and \code{sumz} (community-level sum($z$) curves)
          back into Python.
    \item Overlay TITAN thresholds on the ZCI~vs.~Pollution Score scatter
          plot:
    \begin{itemize}[nosep]
        \item Vertical lines at the Pollution Score community thresholds
              (sum($z^{-}$) peak, sum($z^{+}$) peak).
        \item Horizontal lines at optionally detected ZCI thresholds if
              TITAN is also run with ZCI as a gradient.
        \item Shaded bands for bootstrap confidence intervals.
    \end{itemize}
\end{enumerate}

\paragraph{Theoretical rationale.}
TITAN adds a \textbf{taxon-resolution layer} that the current framework
lacks.  While the ZCI gives a single aggregate condition score per site,
TITAN identifies \emph{which} taxa are responsible for community shifts
and at \emph{what gradient value} they change.  Specifically:
\begin{itemize}[nosep]
    \item It answers whether the \textbf{reference-site threshold} used
          to define $\mathcal{R}$ (bottom 20\,\%) corresponds to an
          ecologically meaningful change point, or is arbitrary.
    \item It reveals whether different taxa respond at different
          pollution levels (gradual turnover) or at a single critical
          threshold (abrupt regime shift).
    \item The community sum($z$) peaks provide an independent estimate
          of the ``reference/degraded boundary'' that can be compared
          with the Stage~5 piecewise-QR breakpoint.
\end{itemize}

\paragraph{Caveats.}
\begin{itemize}[nosep]
    \item TITAN was designed for site $\times$ gradient threshold
          detection; using ZCI as a secondary gradient is a non-standard
          but creative extension.
    \item Results depend on the cluster assignment from Stage~3.
    \item The 5th-to-95th percentile trimming used by TITAN2 means
          extreme-gradient thresholds are not detectable.
    \item Bootstrap resampling ($\ge$500 replicates) is recommended for
          publication; the notebook prototype used 100 for speed.
\end{itemize}


% ═════════════════════════════════════════════════════════════════════════════
\section{Summary of Recommendations}
\label{sec:recommendations}
% ═════════════════════════════════════════════════════════════════════════════

\begin{enumerate}
    \item \textbf{ANOVA assumptions}: add Shapiro--Wilk and
          Levene's tests; report Welch or Kruskal--Wallis as robust
          alternatives where assumptions fail.
    \item \textbf{LDA CV leakage}: move \code{StandardScaler} fitting
          inside each CV fold.
    \item \textbf{NMDS axis-flip caveat}: state explicitly in the thesis
          that NMDS1 orientation is enforced, not discovered.
    \item \textbf{Endpoint sensitivity}: report ZCI stability under
          varying $n_{\mathrm{ep}}$.
    \item \textbf{PCA variance scaling}: note explicitly that PCA is
          performed on log-transformed but not unit-scaled variables, so
          higher-variance contaminants dominate.
    \item \textbf{TITAN integration}: embed per-cluster TITAN into the
          framework after Stage~3, using the Pollution Score gradient and
          relative-abundance taxa matrix; overlay community thresholds on
          ZCI--Pollution plots.
\end{enumerate}


% ═════════════════════════════════════════════════════════════════════════════
% References
% ═════════════════════════════════════════════════════════════════════════════
\begin{thebibliography}{9}
\bibitem{baker2010}
Baker, M.E.\ \& King, R.S.\ (2010).
A new method for detecting and interpreting biodiversity and ecological
community thresholds.
\emph{Methods in Ecology and Evolution}, 1, 25--37.

\bibitem{dufrene1997}
Dufr\^{e}ne, M.\ \& Legendre, P.\ (1997).
Species assemblages and indicator species: the need for a flexible
asymmetrical approach.
\emph{Ecological Monographs}, 67, 345--366.

\bibitem{feng2011}
Feng, X., He, X.\ \& Hu, J.\ (2011).
Wild bootstrap for quantile regression.
\emph{Biometrika}, 98, 995--999.
\end{thebibliography}

\end{document}
